% IEEEtran journal draft. Numbers come from paper/tables/*.tex (paper/make_tables.py) and paper/tables/numbers.json.
% TODO markers must all be resolved before submission. Compile in Overleaf (IEEEtran class).
\documentclass[journal]{IEEEtran}
\usepackage{booktabs,hyperref,cite}
\begin{document}
\title{Does Removing the GIL Raise the Ceiling? A Repeated Cloud Benchmark of Python, Java and Go Services That Wait on Partners}
\author{TODO author names and affiliations}
\maketitle

\begin{abstract}
Integration services such as hotel channel-manager connectors spend most of their time waiting on partner systems and occasionally do CPU work. We compare Go, Java and several Python stacks on the same synthetic connector workload, running every configuration three times on a free shared cloud runner inside containers limited to 2 CPUs and 512~MiB. A cell passes only if all three repeats meet a latency and error SLO. On waiting calls, Go and Java virtual threads sustain 1000 requests in flight, while a 200-thread Java pool fails every request. At 1000 requests per second, median memory over ten repeats ranges from 24~MiB (Go) to 213~MiB (Java virtual threads). For CPU-bound Python, removing the GIL with the free-threaded CPython 3.14 build did not raise the throughput ceiling of an asyncio service with a thread pool, although CPU use spread from at most one core to about two. Running hashing inline on the event loop stops passing at about 15 requests per second, while a four-thread pool passes up to 25. On one runner CPU, the FastAPI service passed at 30 requests per second and failed at 40, matching a one-core limit at about 29~ms of CPU per request. Jobs of the same workflow landed on four different CPU models, which limits cross-job comparison and which we treat as a threat to validity.
\end{abstract}

\begin{IEEEkeywords}
benchmarking, Python, GIL, free-threading, virtual threads, Go, tail latency, cloud performance variability
\end{IEEEkeywords}

\section{Introduction}
Services that connect a property-management system to partners (channel managers, payment providers, loyalty systems) are dominated by waiting: each request makes one or more calls to a slow outside system and returns. The runtime choice affects how many such calls can be in flight, how much memory that costs, and what happens when a request needs real computation. Two recent changes make this choice worth re-testing: Java 21 finalised virtual threads~\cite{jep444}, and CPython 3.13 and 3.14 made a build without the Global Interpreter Lock (GIL) available and then officially supported~\cite{pep703,pep779}.

We ask four questions on a controlled, repeated benchmark: (Q1) how do runtimes differ when many calls wait at once; (Q2) what does each cost in memory at a fixed request rate; (Q3) does the free-threaded Python build raise the ceiling for CPU-bound work; and (Q4) does a faster event loop (uvloop) or a thread pool change the Python result. Our contributions are an open, reproducible harness and workflow; a repeated measurement of these questions on GitHub-hosted runners; and a documented account of where such measurements break down, in particular CPU-model variation between jobs.

\section{Background and Related Work}
\textbf{Free-threaded CPython.} PEP~703 proposes a build option that removes the GIL~\cite{pep703}, and PEP~779 sets the criteria under which that build becomes officially supported but optional~\cite{pep779}. Extensions must declare compatibility, otherwise the interpreter can re-enable the GIL; our images avoid compiled event-loop extensions in the main experiment for this reason. \textbf{Virtual threads.} JEP~444 delivers lightweight threads in JDK~21 that let blocking code scale to many concurrent waits~\cite{jep444}. \textbf{Load models.} Open and closed workload generators produce different results at the same nominal load~\cite{schroeder2006open}; we use an open model (constant arrival rate) so that slow responses do not reduce the offered load~\cite{k6}. \textbf{Tail latency.} Percentile latency, not the mean, determines service quality at scale~\cite{dean2013tail}, so our pass/fail rule is a p99 bound. \textbf{Cloud variability.} Performance measured on public cloud instances varies between instance types and over time, with coefficients of variation from near zero to over 100\% for some microbenchmarks~\cite{laaber2019microbenchmarking}; we therefore repeat every cell and record the CPU model of every job. TODO: add and verify further related work on language and framework benchmarks.

\section{Method}
\subsection{Workload and candidates}
Each service exposes endpoints that mimic connector work: \texttt{/partner} (one call to a fake partner that answers after 100~ms), \texttt{/park} (a call held for about one second, so the request rate equals the number of calls in flight), and \texttt{/pycpu} (about 35~ms of pure-Python arithmetic on the reference CPU, run in a 4-thread pool). Candidates are Go, Java (virtual threads and a 200-thread platform pool), Python with aiohttp, FastAPI and Django on uvicorn, and, for the GIL study, CPython 3.14.3 and the free-threaded 3.14.3t with the GIL off (\texttt{PYTHON\_GIL=0}), plus a control that runs the free-threaded build with the GIL forced on.
\subsection{Execution}
Each run starts a fresh set of containers: the candidate limited to 2 CPUs and 512~MiB without swap, a fake partner, the k6 load generator~\cite{k6} and a resource sampler reading cgroup counters once per second from outside the candidate. A run has a warm-up (2~min, discarded), a measured phase (3~min) and a recovery phase (60~s). The candidate is pinned to CPUs 0--1 of a 4-vCPU GitHub-hosted Ubuntu 24.04 runner, and the other components to CPUs 2--3.
\subsection{Pass/fail and repeats}
A run passes if its failure rate is at most $10^{-3}$, no iterations were dropped, and p99 is within the endpoint SLO (250~ms for \texttt{/partner}, 1250~ms for \texttt{/park}, 2000~ms for \texttt{/pycpu}; thresholds fixed in the repository before the first run). Each cell is repeated three times and we report ``k/n'': passing repeats over valid repeats. A cell is called passing only if all valid repeats pass. With three repeats we do not compute confidence intervals; we give exact counts and min--max ranges, and we do not claim statistical significance. CPU cores used are estimated as rate times median CPU time per successful request.
\subsection{Reproducibility}
All code, the workflow and the raw artifacts are in the repository, and each result cites its workflow run. TODO add repository URL and an archive DOI (for example a Zenodo release) before submission.

\section{Results}
\textbf{Data.} The base campaign comprises 129 valid runs (workflow runs \#2--\#4). Two follow-ups address uvloop, hashing and sensitivity questions: run \#5 has 95 valid runs of 96 and run \#6 has 96 runs, of which 8 are invalid because no replies were received (all Django on waiting calls and the 2000-thread Java pool at 1000 in flight).
\subsection{Waiting calls}
\begin{table}[t]
\caption{Waiting calls (\texttt{/park}, each request held about 1\,s by the partner). Cells: repeats passing / valid repeats, median p99, CPU cores used (passing cells). SLO: p99 $\le$ 1250\,ms.}
\label{tab:park}
\centering\small
\begin{tabular}{lcc}
\toprule
Candidate & 250 in flight & 1000 in flight \\
\midrule
Go & 3/3, 1002\,ms, 0.08\,c & 3/3, 1002\,ms, 0.27\,c \\
Java (virtual threads) & 3/3, 1003\,ms, 0.10\,c & 3/3, 1112\,ms, 0.48\,c \\
Java (200 platform threads) & 0/3 (no replies) & 0/3 (no replies) \\
\bottomrule
\end{tabular}
\end{table}

Go and Java virtual threads passed 250 and 1000 calls in flight in every repeat (Table~\ref{tab:park}). The 200-thread Java pool returned no usable replies at either rate in any repeat, which is consistent with the pool being exhausted when the number of concurrent waits exceeds its size. We did not instrument the pool to confirm this.
\subsection{Memory at 1000 requests per second}
\begin{table}[t]
\caption{One partner call per request at 1000\,RPS in a 512\,MiB container. Memory is the median over the measured phase, shown as median [min--max] across repeats.}
\label{tab:memory}
\centering\small
\begin{tabular}{lcccc}
\toprule
Candidate & Pass & p99 (ms) & Memory (MiB) & CPU (ms/req) \\
\midrule
Go & 3/3 & 102 & 23 [23--23] & 0.26 \\
Python aiohttp & 3/3 & 102 & 36 [36--36] & 0.32 \\
FastAPI (uvloop, inline) & 3/3 & 102 & 55 [55--55] & 0.40 \\
Java (virtual threads) & 3/3 & 112 & 217 [213--217] & 0.34 \\
\bottomrule
\end{tabular}
\end{table}

All four candidates met the SLO in every repeat (Table~\ref{tab:memory}). Median memory ranged from 23~MiB (Go) to 217~MiB (Java virtual threads), roughly nine times higher, with CPU time per request in a narrow band of 0.26--0.40~ms. We did not tune the JVM heap, so the Java figure reflects default behaviour in a 512~MiB container, not a minimum.
\subsection{GIL ceiling and forced-GIL control}
\begin{table}[t]
\caption{Pure-Python CPU work (\texttt{/pycpu}), one event loop plus a 4-thread pool. Cells as in Table~\ref{tab:park}; SLO p99 $\le$ 2000\,ms.}
\label{tab:gil}
\centering\small
\begin{tabular}{lccccc}
\toprule
Candidate & 10 RPS & 15 RPS & 20 RPS & 30 RPS & 40 RPS \\
\midrule
aiohttp, GIL & 3/3, 35\,ms, 0.34\,c & 3/3, 35\,ms, 0.51\,c & 3/3, 35\,ms, 0.67\,c & 3/3, 31\,ms, 0.86\,c & 0/3 (no replies) \\
aiohttp, 3.14t with GIL forced on & 3/3, 35\,ms, 0.35\,c & 3/3, 36\,ms, 0.53\,c & 3/3, 35\,ms, 0.70\,c & 3/3, 33\,ms, 0.87\,c & 0/3 (no replies) \\
aiohttp, 3.14t GIL off & 3/3, 35\,ms, 0.35\,c & 3/3, 36\,ms, 0.52\,c & 3/3, 36\,ms, 0.70\,c & 3/3, 46\,ms, 1.26\,c & 0/3 (no replies) \\
\bottomrule
\end{tabular}
\end{table}

All three builds passed up to 30~RPS and none passed at 40~RPS (Table~\ref{tab:gil}). Under the GIL, CPU use stayed below one core at every passing rate, as expected, while the free-threaded build used about 1.3 cores at 30~RPS and about two at 40~RPS. Despite this, the free-threaded build did not pass at a higher rate. We caution that the 30~RPS rows are not strictly comparable: the GIL builds' 30 and 40~RPS runs executed on an AMD EPYC 9V74 runner and the free-threaded build's on an EPYC 7763; where all builds shared one CPU model (20~RPS) CPU per request agreed within about 3\%. The failures at 40~RPS were client timeouts with essentially no successful replies; we infer overload collapse from this but have no server-side evidence, because candidate logs were empty.
\begin{table}[t]
\caption{FastAPI and Django on \texttt{/pycpu} with plain asyncio and a 4-thread pool (GIL build).}
\label{tab:frameworks}
\centering\small
\begin{tabular}{lccc}
\toprule
Candidate & 10 RPS & 15 RPS & 20 RPS \\
\midrule
FastAPI, GIL, asyncio, 4 threads & 3/3, 37\,ms, 0.36\,c & 3/3, 38\,ms, 0.54\,c & 3/3, 38\,ms, 0.71\,c \\
Django, GIL, asyncio, 4 threads & 3/3, 41\,ms, 0.36\,c & 3/3, 38\,ms, 0.54\,c & 3/3, 38\,ms, 0.71\,c \\
\bottomrule
\end{tabular}
\end{table}

\subsection{Effect of uvloop and a thread pool}
\begin{table}[t]
\caption{Effect of uvloop on \texttt{/pycpu} (GIL build, 4-thread pool). FastAPI and Django jobs ran on different runner CPUs, see Section~\ref{sec:threats}.}
\label{tab:uvloop}
\centering\small
\begin{tabular}{lccccc}
\toprule
Candidate & 10 & 15 & 20 & 30 & 40 \\
\midrule
FastAPI, GIL, asyncio, 4 threads & -- & -- & -- & 3/3, 23\,ms, 0.67\,c & 3/3, 23\,ms, 0.89\,c \\
FastAPI, GIL, uvloop, 4 threads & 3/3, 23\,ms, 0.22\,c & 3/3, 23\,ms, 0.33\,c & 3/3, 23\,ms, 0.44\,c & 3/3, 23\,ms, 0.66\,c & 3/3, 23\,ms, 0.88\,c \\
Django, GIL, asyncio, 4 threads & -- & -- & -- & 0/3 (no replies) & 0/3 (no replies) \\
Django, GIL, uvloop, 4 threads & 3/3, 37\,ms, 0.36\,c & 3/3, 39\,ms, 0.54\,c & 3/3, 39\,ms, 0.71\,c & 0/3 (no replies) & 0/3 (no replies) \\
\bottomrule
\end{tabular}
\end{table}

Replacing the asyncio loop by uvloop did not change the ceiling in either framework (Table~\ref{tab:uvloop}). For FastAPI both variants passed up to 40~RPS with p99 near 23~ms, but this job ran on an Intel Xeon 6975P-C where the same work needs about 22~ms of CPU instead of about 36~ms; its ceiling was not reached. Django, on the AMD runner, passed at 20~RPS and failed at 30~RPS in both variants. Section~\ref{sec:followup} gives the remaining checks.

\subsection{Hashing, higher rates and waiting-call sensitivity}\label{sec:followup}
\begin{table}[t]
\caption{Hashing (\texttt{/cpu}, releases the GIL), uvloop in all rows. The FastAPI rows ran on an AMD EPYC 9V74 runner and the Django rows on an EPYC 7763 runner, so compare within a framework only. SLO: p99 $\le$ 2000\,ms.}
\label{tab:hashing}
\centering\small
\begin{tabular}{lcccc}
\toprule
Candidate & 10 RPS & 15 RPS & 20 RPS & 25 RPS \\
\midrule
FastAPI, inline on the event loop & 3/3, 59\,ms, 0.59\,c & 3/3, 61\,ms, 0.88\,c & 0/3 (no replies) & 0/3 (no replies) \\
FastAPI, 4-thread pool & 3/3, 60\,ms, 0.59\,c & 3/3, 63\,ms, 0.89\,c & 3/3, 65\,ms, 1.20\,c & 3/3, 65\,ms, 1.54\,c \\
Django, inline on the event loop & 3/3, 68\,ms, 0.68\,c & 0/3, 2639\,ms & 0/3 (no replies) & 0/3 (no replies) \\
Django, 4-thread pool & 3/3, 71\,ms, 0.68\,c & 3/3, 71\,ms, 1.02\,c & 3/3, 73\,ms, 1.38\,c & 3/3, 74\,ms, 1.74\,c \\
\bottomrule
\end{tabular}
\end{table}

For hashing, which releases the GIL, running the work inline on the event loop stopped passing at 20~RPS for FastAPI and at 15~RPS for Django (Table~\ref{tab:hashing}), consistent with one core divided by about 60--70~ms of CPU per request. With a four-thread pool both passed every repeat up to 25~RPS using up to 1.7 cores; run \#5 saw failures at 30~RPS, so the pooled ceiling lies between 25 and 30~RPS, consistent with two cores. This is the clearest effect of concurrency design in our data.
\begin{table}[t]
\caption{FastAPI on \texttt{/pycpu} above 40\,RPS (AMD EPYC 7763 runner, about 36\,ms of CPU per request).}
\label{tab:pycpuhigh}
\centering\small
\begin{tabular}{lccc}
\toprule
Candidate & 50 RPS & 60 RPS & 80 RPS \\
\midrule
FastAPI, GIL, asyncio, 4 threads & 0/3 (no replies) & 0/3 (no replies) & 0/3 (no replies) \\
FastAPI, GIL, uvloop, 4 threads & 0/3 (no replies) & 0/3 (no replies) & 0/3 (no replies) \\
\bottomrule
\end{tabular}
\end{table}

For pure-Python work, FastAPI failed every request at 50, 60 and 80~RPS on an AMD EPYC 7763 runner (Table~\ref{tab:pycpuhigh}), where one request needs about 36~ms of CPU, so 50~RPS needs about 1.8 cores. It had passed 40~RPS on an Intel runner at about 22~ms (0.88 core). Because those two sides ran on different CPUs, we repeated the bracket on a single runner CPU in run \#7 (Table~\ref{tab:bracket}).
\begin{table}[t]
\caption{FastAPI on \texttt{/pycpu} bracketed on a single runner CPU (Intel Xeon Platinum 8573C, both builds in one job). About 29\,ms of CPU per request, so one core allows about 34 RPS.}
\label{tab:bracket}
\centering\small
\begin{tabular}{lcccc}
\toprule
Candidate & 30 RPS & 40 RPS & 50 RPS & 60 RPS \\
\midrule
FastAPI, GIL, asyncio, 4 threads & 3/3, 30\,ms, 0.86\,c & 0/3 (no replies) & 0/3 (no replies) & 0/3 (no replies) \\
FastAPI, GIL, uvloop, 4 threads & 3/3, 30\,ms, 0.85\,c & 0/3 (no replies) & 0/3 (no replies) & 0/3 (no replies) \\
\bottomrule
\end{tabular}
\end{table}

On an Intel Xeon Platinum 8573C runner both FastAPI builds passed 30~RPS in every repeat at about 29~ms of CPU per request (0.86 core) and returned no usable replies at 40, 50 and 60~RPS. One core allows 1000/29, about 34~RPS, so the ceiling lies between 30 and 40~RPS on one CPU model and is consistent with a one-core limit. uvloop made no difference. The 40~RPS failure is well above the limit (1.16 cores needed), so we do not claim the exact 34~RPS point.
\subsection{Ten repeats of the headline cells}\label{sec:repeats}
\begin{table}[t]
\caption{Waiting calls with 10 repeats per cell (invalid repeats are excluded from n). Pass column: passing/valid repeats with the exact 95\% Clopper--Pearson interval for the pass probability. p99 and memory are median [min--max].}
\label{tab:repeatspark}
\centering\small
\begin{tabular}{lcccc}
\toprule
Candidate & In flight & Pass [95\% CI] & p99 (ms) & Memory (MiB) \\
\midrule
Go & 250 & 10/10 [69--100\%] & 1003 [1003--1003] & 35 [33--35] \\
Go & 1000 & 10/10 [69--100\%] & 1002 [1002--1002] & 116 [116--116] \\
Java (virtual threads) & 250 & 10/10 [69--100\%] & 1005 [1004--1006] & 215 [212--218] \\
Java (virtual threads) & 1000 & 10/10 [69--100\%] & 1137 [1083--1176] & 229 [225--232] \\
Java (200 platform threads) & 250 & 0/10 [0--31\%] & -- & 247 [244--266] \\
Java (200 platform threads) & 1000 & 0/10 [0--31\%] & -- & 322 [293--356] \\
Java (2000 platform threads) & 250 & 10/10 [69--100\%] & 1013 [1011--1016] & 481 [459--486] \\
Java (2000 platform threads) & 1000 & 0/6 [0--46\%] & -- & 506 [495--510] \\
\bottomrule
\end{tabular}
\end{table}

Run \#7 repeated the waiting-call and memory cells ten times (Tables~\ref{tab:repeatspark} and~\ref{tab:repeatsmem}). Go and Java virtual threads passed all ten repeats at 250 and 1000 calls in flight; with ten of ten, the exact 95\% lower bound on the pass probability is 69\%, so we do not claim a pass rate above that. The 200-thread Java pool failed all ten repeats at both rates (upper bound 31\%). The 2000-thread pool passed 10 of 10 at 250; at 1000 in flight six repeats were valid and failed, and four were invalid because the candidate returned no replies at all, so we count the cell as failing. The park cells ran on different CPU models (Go and virtual-thread Java on an AMD EPYC 7763, the platform pools on an EPYC 9V74), which is a limitation, but the outcomes are all-or-nothing rather than marginal.
\begin{table}[t]
\caption{One partner call per request at 1000\,RPS, 10 repeats per candidate. Memory and p99 are median [min--max] over repeats.}
\label{tab:repeatsmem}
\centering\small
\begin{tabular}{lccc}
\toprule
Candidate & Pass [95\% CI] & Memory (MiB) & p99 (ms) \\
\midrule
Go & 10/10 [69--100\%] & 24 [23--24] & 102 [102--102] \\
Python aiohttp & 10/10 [69--100\%] & 36 [36--36] & 102 [102--102] \\
FastAPI (uvloop, inline) & 10/10 [69--100\%] & 55 [55--55] & 102 [102--102] \\
Java (virtual threads) & 10/10 [69--100\%] & 213 [204--217] & 111 [110--111] \\
\bottomrule
\end{tabular}
\end{table}

Memory was stable across repeats: Go 24~MiB [23--24], aiohttp 36 [36--36], FastAPI 55 [55--55] and Java virtual threads 213 [204--217], so Java used about nine times the memory of Go. All four candidates ran on one Intel Xeon Platinum 8573C job.
\begin{table}[t]
\caption{Waiting-call sensitivity runs on an Intel Xeon 6973P-C runner. Columns: \texttt{/partner} at 1000\,RPS, then \texttt{/park} with 250 and 1000 in flight; last column is median memory in MiB for the three. A cell without replies means the candidate returned no usable answers.}
\label{tab:waiting2}
\centering\small
\begin{tabular}{lcccc}
\toprule
Candidate & partner 1000 & park 250 & park 1000 & Memory (MiB) \\
\midrule
FastAPI, GIL, asyncio, 4 threads & 3/3, 102\,ms, 0.33\,c & 3/3, 1003\,ms, 0.10\,c & 3/3, 1003\,ms, 0.37\,c & 71 / 75 / 113 \\
Java (200 platform threads) & 3/3, 112\,ms, 0.19\,c & -- & -- & 242 / -- / -- \\
Java (2000 platform threads) & -- & 3/3, 1016\,ms, 0.09\,c & 0/2 (no replies) & -- / 469 / 501 \\
\bottomrule
\end{tabular}
\end{table}

On waiting calls, FastAPI passed every repeat with both asyncio and uvloop, at 70--113~MiB (Table~\ref{tab:waiting2}). A Java pool of 2000 platform threads passed 250 calls in flight but returned no replies at 1000, with memory at 501~MiB of the 512~MiB limit; a pool of 200 threads passed 1000 requests per second of 100~ms calls (242~MiB), because each call is short. Pool size therefore matters only when calls are held long enough to exhaust it.

\section{Threats to Validity}\label{sec:threats}
\textbf{Hardware variation.} GitHub does not let us choose the runner CPU. Jobs in this study ran on AMD EPYC 7763, AMD EPYC 9V74, Intel Xeon 6975P-C and Intel Xeon 6973P-C hosts; we record the model per job and avoid cross-job comparisons where it differs. \textbf{Shared runners.} Co-tenant noise is possible. Ten repeats of the same cell on one job gave memory ranges from 0\% (aiohttp, FastAPI) to about 6\% (Java virtual threads) and p99 ranges under 1\% except Java at 1000 calls in flight (1083--1176~ms); we did not repeat cells across jobs, where CPU-model differences dominate. \textbf{Scope.} Three repeats and 3-minute measured phases; one synthetic workload; one runner family; the Java heap was not tuned; free-threaded results depend on the interpreter version and on libraries that do not re-enable the GIL. \textbf{Django on waiting calls.} A Django test application failed waiting-call tests (about 30\% failures at 250 calls in flight, 400~MiB) while FastAPI passed. The resource samples show one operating-system thread per request in flight (about 254 threads at 250 in flight, against 3 for FastAPI). Reading Django 5.2 and asgiref 3.13 shows the cause we believe: each request runs inside a thread-sensitive context and the default \texttt{request\_started} and \texttt{request\_finished} receivers are synchronous, so each request gets its own thread. A local control without resource limits showed 252 threads falling to 1 when those receivers were removed. TODO: replace this sentence with the controlled benchmark result of batch \texttt{django-signals}. \textbf{Causal claims.} We report outcomes and resource use. We do not profile the interpreter, so explanations such as the GIL cap or pool exhaustion are inferences from rates and CPU use, not measurements of the mechanism.

\section{Conclusion}
Our data support four findings on a 2-CPU, 512~MiB budget with a 100~ms partner: (1) virtual threads and goroutines sustain 1000 concurrent one-second waits, while a 200-thread pool does not, and a 2000-thread pool fails at 1000 near the memory cap; (2) memory at a fixed rate differs by about an order of magnitude between Go and Java virtual threads; (3) on this workload the free-threaded build spreads CPU use across cores without raising the passing ceiling, and neither uvloop nor the thread pool moved the pure-Python ceiling; (4) for GIL-releasing work, a thread pool moves the ceiling from about one core to about two; and (5) on one runner CPU the FastAPI pure-Python ceiling lies between 30 and 40~RPS, consistent with one core. The one open item is the controlled confirmation of the Django thread-per-request explanation.

\section*{Acknowledgment}
TODO: This work used an AI assistant (Claude, Anthropic) to help write the benchmark tooling, analyse the data and draft the text. The authors reviewed the content and are responsible for it. Confirm the wording against the target publisher's AI-disclosure policy.

\bibliographystyle{IEEEtran}
\bibliography{refs}
\end{document}
